
        \documentclass[fleqn]{article}      
        \usepackage{amsmath}
        \usepackage{fontspec}
        \usepackage{kotex} % 한국어 지원  

        \begin{document}      
        \section*{쉬운 문제}

\noindent 문제 1. 다음 행렬 \( A = \begin{pmatrix} 1 & 2 \\ 3 & 4 \end{pmatrix} \)의 주 대각선의 합을 구하시오. \\\\

\noindent 문제 2. 두 행렬 \( A = \begin{pmatrix} 1 & 0 \\ 0 & 1 \end{pmatrix} \)와 \( B = \begin{pmatrix} 2 & 3 \\ 4 & 5 \end{pmatrix} \)의 합을 구하시오. \\\\

\noindent 문제 3. 다음 행렬의 전치 행렬을 구하시오: \( A = \begin{pmatrix} 2 & 4 \\ 5 & 7 \end{pmatrix} \) \\\\

\noindent 문제 4. 벡터 \( \begin{pmatrix} 3 \\ 2 \end{pmatrix} \)와 스칼라 \( 2 \)의 곱을 구하시오. \\\\

\noindent 문제 5. 행렬 \( A = \begin{pmatrix} 1 & 2 \\ 3 & 4 \end{pmatrix} \)의 행렬식을 구하시오. \\\\

\section*{중간 문제}

\noindent 문제 6. 다음 행렬의 곱을 구하시오: \( A = \begin{pmatrix} 1 & 2 \\ 3 & 4 \end{pmatrix} \)와 \( B = \begin{pmatrix} 5 & 6 \\ 7 & 8 \end{pmatrix} \) \\\\

\noindent 문제 7. 행렬 \( A = \begin{pmatrix} 2 & 3 \\ 4 & 5 \end{pmatrix} \)의 역행렬을 구하시오. \\\\

\noindent 문제 8. 다음 두 행렬 \( A = \begin{pmatrix} 1 & 2 \\ 3 & 4 \end{pmatrix} \)와 \( B = \begin{pmatrix} 5 & 6 \\ 7 & 8 \end{pmatrix} \)의 차를 구하시오. \\\\

\noindent 문제 9. \( A = \begin{pmatrix} 1 & 3 \\ 0 & 2 \end{pmatrix} \)와 \( B = \begin{pmatrix} 4 & 1 \\ 2 & 0 \end{pmatrix} \)의 곱을 구하시오. \\\\

\noindent 문제 10. 행렬 \( A = \begin{pmatrix} 2 & -1 \\ 1 & 3 \end{pmatrix} \)의 고유값을 구하시오. \\\\

\section*{어려운 문제}

\noindent 문제 11. 행렬 \( A = \begin{pmatrix} 1 & 1 & 1 \\ 2 & 2 & 2 \\ 3 & 3 & 3 \end{pmatrix} \)의 랭크를 구하시오. \\\\

\noindent 문제 12. \( A = \begin{pmatrix} 2 & 3 \\ 1 & 4 \end{pmatrix} \)와 \( B = \begin{pmatrix} 5 & 6 \\ 7 & 8 \end{pmatrix} \)의 곱의 고유값을 구하시오. \\\\

\noindent 문제 13. 행렬 \( A = \begin{pmatrix} 2 & 0 & 1 \\ 3 & 1 & 0 \\ 0 & 4 & 2 \end{pmatrix} \)의 행렬식을 구하시오. \\\\

\noindent 문제 14. 다음 행렬의 특성 방정식을 구하시오: \( A = \begin{pmatrix} 5 & 4 \\ 2 & 3 \end{pmatrix} \) \\\\

\noindent 문제 15. \( A = \begin{pmatrix} 1 & 2 \\ 2 & 4 \end{pmatrix} \)와 \( B = \begin{pmatrix} 1 & 0 \\ 0 & 1 \end{pmatrix} \)의 곱의 랭크를 구하시오. \\\\

\section*{답안}

1. 5 \\\\
2. \( \begin{pmatrix} 3 & 3 \\ 4 & 6 \end{pmatrix} \) \\\\
3. \( A^T = \begin{pmatrix} 2 & 5 \\ 4 & 7 \end{pmatrix} \) \\\\
4. \( \begin{pmatrix} 6 \\ 4 \end{pmatrix} \) \\\\
5. -2 \\\\
6. \( \begin{pmatrix} 19 & 22 \\ 43 & 50 \end{pmatrix} \) \\\\
7. \( \frac{1}{-1} \begin{pmatrix} 5 & -3 \\ -4 & 2 \end{pmatrix} = \begin{pmatrix} -5 & 3 \\ 4 & -2 \end{pmatrix} \) \\\\
8. \( \begin{pmatrix} -4 & -4 \\ -4 & -4 \end{pmatrix} \) \\\\
9. \( \begin{pmatrix} 14 & 3 \\ 14 & 0 \end{pmatrix} \) \\\\
10. 1, 4 \\\\
11. 1 \\\\
12. 0, 10 \\\\
13. -12 \\\\
14. \( \lambda^2 - 8\lambda + 7 = 0 \) \\\\
15. 1 \\\\      
        \end{document} 
    